\section{Experimental Results}
\label{sec:results}

Literature on benchmarks for databases and data store is extensive
and covers what measures to collect for various guarantee of consistency levels~\cite{ConsistencyOverview}.
As our database has no consistency guarantees, we focus our experimentation on latency and scalability.

Experimental data were collected using three machines:
\begin{itemize}
    \item Machine $M$: MacBook Pro with M1 and \qty{32}{\giga\byte} RAM, built-in storage, operating system macOS~26, file system \texttt{APFS};
    \item Machine $A$: Intel Core i7-8750H (6 cores, 12 threads with Hyper Threading, locked at \qty{2.2}{\giga\hertz}), \qty{16}{\giga\byte} DDR4 RAM, \qty{128}{\giga\byte} SATA3 SSD, operating system Arch Linux, kernel \texttt{6.17.1}, file system \texttt{btrfs};
    \item Machine $H$: Intel Core i5-12450H (4 E-cores and 4 P-Cores, 12 threads with Hyper Threading, E-Cores locked at \qty{1.5}{\giga\hertz} and P-Cores at \qty{2}{\giga\hertz}), \qty{16}{\giga\byte} DDR4 RAM, \qty{512}{\giga\byte} NVMe 4xPCIe3 SSD, operating system Arch Linux, kernel \texttt{6.17.1}, file system \texttt{XFS}.
\end{itemize}

The machines were connected via a \qty{1}{\giga\bit} Ethernet network with a latency of about \qty{0.3}{\milli\second} while under stress.
For these tests, we configured the system with a \qty{4}{\kilo\byte} max request size,
a value chosen to avoid saturating the network connection in most tests,
and which allowed us to measure the performance of CPU and disk.

A purpose-built program (\texttt{ktb}) was developed to stress the system using multiple connections.
We first recorded baseline values by having each machine stress itself in both \texttt{item-get} and \texttt{item-set} operations
on a table with throttling disabled (i.e., set to sufficiently high limit),
1000 keys and value sizes of \qty{4}{\kilo\byte}.

Table~\ref{tab:results-stress} presents the results, across the \texttt{item-get} at \qty{4}{\kilo\byte} test, we observe $H > M > A$;
while for \texttt{item-set} at \qty{4}{\kilo\byte}, $2H \approx M \approx A$.
Whenever the \texttt{item-get} at \qty{4}{\kilo\byte} run saturated the \qty{1}{\giga\bit} connection,
we also report the operation at \qty{512}{\byte}.

The scalability and replication of the system were then evaluated.
Unless stated otherwise, the following tests were conducted by running \texttt{ktb} on machine $H$, considering only \texttt{item-set} operations.
\texttt{item-get} operations were excluded because the operating system caches file reads, making their results unrepresentative of true random-access disk performance.
Nevertheless, the measured req/s values for \texttt{item-get} operations are reported in Table~\ref{tab:results-stress}.

We configured the machines for either a uniform key distribution, where each machine's bandwidth parameter was set to the same sufficiently high value,
or with a proportional key distribution, where each machine's bandwidth was still set sufficiently high but proportional to its self-stress results.
In our setup, for simplicity, machines $M$ and $A$ were assigned the same value, while $H$ was configured with twice that amount.

\paragraph{Statistical Tests}
As introduced in Section~\ref{subsubsec:stddev} the following tests
may measure a slowdown due to the statistical nature of key distribution.

For each test and for each node we independently assess whether the measured slowdown $f$ is attributable to the key distribution by:
\begin{itemize}
    \item Computing the expected $\delta_i$ using Equation~\ref{eq:slowdown}:
    \[
        \delta_i = (f - 1) \cdot \sqrt{\frac{|K| \cdot p_i}{1 - p_i}}
    \]
    \item Testing the null hypothesis:
    \begin{center}
        \textit{For each node, keys are assigned according to a Binomial distribution with $N=1000$ and probability $p_i$ as specified by the configuration.}
    \end{center}
    We compute the one-sided $p$-value using a Gaussian approximation.
\end{itemize}

Rejecting the null hypothesis implies the presence of an additional source of slowdown
in the tested configuration, most likely CPU bottlenecks.
If we fail to reject the null hypothesis,
the observed performance does not provide sufficient evidence against it,
so the measured slowdown is \textit{likely} due to statistical variability of key assignment.

%def make_numbers(actual, expected, prop):
%    f = actual / expected
%    f = 1 / f
%    d = (f - 1) * sqrt(1000 * prop / (1 - prop))
%    p = Fraction(1, 2) * erfc(d / sqrt(2))
%    print(f'{d.n()=}, {p.n()=}')

\paragraph{Distribution}
To test key distribution and scalability we used a table configured
with sufficient bandwidth to saturate the test case, $N=1$, $R=1$, and $W=1$.
This configuration distributes requests across peers according to each machine's bandwidth proportion.

In the uniform case each machine receives one third of the requests,
so the slowest machine becomes the bottleneck.
Therefore, the overall throughput should be approximately three times the slowest machine's throughput.
In our setup the slowest machine is $A$, which achieved \qty{4.8}{\kilo\hertz},
yielding an expected overall throughput of \qty{14.4}{\kilo\hertz}.
During testing the system reached \qty{13.2}{\kilo\hertz} (\qty{92}{\percent} of expected),
producing a delta of \num{2.0328} and a $p$-value of \num{0.0210}.
This $p$-value gives marginal evidence against the null hypothesis;
thus the slowdown is likely due to both a slightly imbalanced distribution and additional CPU overhead.

% p_1,2,3 = 1/3
% make_numbers(Fraction(132, 10), Fraction(144, 10), 1/3)
% d.n()=2.03278907045435, p.n()=0.0210369189124584

In the proportional case, as long as the machine bandwidth are configured correctly no single machine should become a bottleneck
and the overall throughput should equal the sum of individual throughputs.
In our setup, this corresponds to \qty{20.6}{\kilo\hertz},
while the test measured \qty{17.0}{\kilo\hertz} or \qty{83}{\percent} of the expected value.
For machines $M$ and $A$ the deltas are \num{3.8663} and the $p$-value is \num{5.5255e-5},
while for machine $H$ the delta is \num{6.6966} and its $p$-value can be safely approximated with zero.
These $p$-values allow us to reject the null hypothesis,
indicating a significant slowdown due to CPU overhead.

% p_1 = p_2 = 1/4, p_3 = 1/2
% make_numbers(Fraction(170, 10), Fraction(206, 10), 1/2)
% d.n()=6.69658798623892, p.n()=1.06671009488812e-11
% make_numbers(Fraction(170, 10), Fraction(206, 10), 1/4)
% d.n()=3.86627687650705, p.n()=5.52547364766706e-5

\paragraph{Replication}
To test replication we used a table with sufficient bandwidth, $N=2$, $R=2$, and $W=2$.
This configuration provides reproducible performance ($N=R=W$) and was preferred over the $N=3$, $R=3$, $W=3$ case.

In the uniform replication case each machine receives two thirds of client requests, as the key space is replicated twice,
so the expected overall throughput is $3/2$ times the slowest machine's throughput.
The expected overall throughput is \qty{7.2}{\kilo\hertz}, and testing confirms the system achieves \qty{7.2}{\kilo\hertz}.

% p_1,2,3 = 2/3
% make_numbers(Fraction(72, 10), Fraction(72, 10), 2/3)
% d.n()=0, p.n()=0.500000000000000

In the proportional replication case node $H$ holds a full copy while $M$ and $A$ each store half.
The overall throughput is therefore the minimum of $H$'s throughput and the sum of $M$ and $A$'s rates.
In our setup $H$ is faster than $M$ and $A$ combined, giving an expected overall throughput of \qty{9.8}{\kilo\hertz}.
The system achieved \qty{9.4}{\kilo\hertz} or \qty{96}{\percent} of the expected value.
For $M$ and $A$ the deltas are \num{1.3457} and the $p$-value is\num{0.0892};
while for $H$ the delta is not significant because it stores exactly a full copy.
Here the $p$-values are sufficiently large that we cannot reject the null hypothesis,
so the observed slowdown is likely due to an imbalanced key distribution.


% p_1,2 = 1/2, p_3=1
% make_numbers(Fraction(94, 10), Fraction(98, 10), 1/2)
% d.n()=1.34565006815676, p.n()=0.0892076982478822
% meaningless otherwise

\begin{table}
    \setlength{\tabcolsep}{3pt}

    \centering
    \caption{Maximum request frequencies}
    \label{tab:results-stress}


    \sisetup{table-format = 2.1{\textsuperscript{\emph{A}}}}
    \begin{tabular}{lSSS}
        \toprule
        \multirow{2}{*}{\raisebox{-0.5ex}{Test Case}} & \multicolumn{3}{c}{Throughput (\unit{\kilo\hertz})} \\
        \cmidrule{2-4}
        & \texttt{get}~\qty{512}{\byte} & \texttt{get}~\qty{4}{\kilo\byte} & \texttt{set}~\qty{4}{\kilo\byte} \\
        \midrule
        Machine $M$ self stress   & \textendash                   & 64.0                             & 5.0                              \\
        Machine $H$ self stress   & \textendash                   & 83.3                             & 10.8                             \\
        Machine $A$ self stress   & \textendash                   & 56.6                             & 4.8                              \\
        Uniform with $N = 1$      & 102.0                         & 44.1                             & 13.2                             \\
        Proportional with $N = 1$ & 97.0                          & 56.5                             & 17.0                             \\
        Uniform with $N = 2$      & 67.3                          & 21.8                             & 7.2                              \\
        Proportional with $N = 2$ & 62.4                          & 28.0                             & 9.4                              \\
        \bottomrule
    \end{tabular}
\end{table}

%\footnotetext{These tests were limited by the \qty{1}{\giga\bit} connection between the machines.}

\paragraph{Real-Time}
The last test focused on evaluating the real-time properties of the system, simulating a concrete use case.
With a proportional bandwidth configuration matching the measured rates in the self stress,
a table $T_T$ was created with $N=3$, $R=3$, $W=3$ and a bandwidth of \qty{200}{\hertz}.
To represent external interference, another table $T_O$ was created with $N=2$, $R=2$, $W=2$
and bandwidth such that it would saturate the machines configured bandwidth.
To stress the system, two extra instances of \texttt{ktb} were executed on machines $A$ and $M$, both targeting $T_O$.
The reported data was collected from machine $H$'s \texttt{ktb} instance targeting $T_T$.

The system successfully isolated the target table $T_T$ from the outside noise, maintaining the allocated throughput of \qty{200}{\hertz}.
Reported in Table~\ref{tab:results-isolation} are the measured latencies when the target table received less than \qty{200}{\hertz},
both with and without external stress.
Overall, a slight increase in latency was observed, which is expected, as no explicit latency control mechanism is implemented.

\begin{table}
    \centering
    \caption{Latency percentiles}
    \label{tab:results-isolation}
    \begin{tabular}{lrrrrrr}
        \toprule
        \multirow{2}{*}{\raisebox{-0.5ex}{Test Case}} & \multicolumn{5}{c}{Latency (\unit{\milli\second})} \\
        \cmidrule{2-6}
        & Average & $P_{50}$ & $P_{90}$ & $P_{99}$ & $P_{99.9}$ \\
        \midrule
        Baseline on H   & 1.28    & 1.22     & 1.38     & 1.63     & 19.50      \\
        Stressed on H   & 2.35    & 2.14     & 2.78     & 6.12     & 37.87      \\
        Dist.\ Baseline & 4.18    & 3.22     & 6.75     & 16.09    & 37.55      \\
        Dist.\ Stressed & 12.10   & 11.85    & 19.37    & 48.29    & 92.33      \\
        \bottomrule
    \end{tabular}
\end{table}

% TODO: Una sezione che spiega due cose in piu` ci potrebbe stare, viene anche decentemente lunga.
%\subsection{Key Distribution Error}
%\label{subsec:key-distribution-error}
